\section{Background}
\label{sec:background}


\noindent \textbf{Blockchain:}
Blockchain is a distributed ledger that broadcasts and stores information of
transactions across different parties. 
A blockchain consists of a growing
number of blocks and a consensus algorithm determining block order.
Each block is constituted of transactions.
Ethereum~\cite{ethereum22,buterin2014ethereum} is the first blockchain to
support, store, and execute Turing complete programs, known as smart contracts.
Many new blockchains use the Ethereum virtual machine (EVM) 
for execution due to its popularity among developers.

%Smart contracts can create internal transactions to other smart contracts. The internal transactions are nested from a top-level transaction, an external transaction, initiated by a user account. If a transaction is aborted then the effects of all nested internal transactions will be reversed. 



\noindent \textbf{Smart Contracts:}
Each smart contract is associated with a unique address, a
persistent account's storage trie, a balance of native tokens, e.g., Ether in
Ethereum, and bytecode (e.g., EVM
bytecode~\cite{ethereum22,buterin2014ethereum}) that executes incoming
transactions to change the storage and balance. Users interact with a smart
contract by issuing transactions from their user accounts to the contract
address. 
Smart contracts can also interact with other smart contracts as function calls.
Currently, there are several human-readable high-level programming languages,
e.g., Solidity~\cite{solidity} and Vyper~\cite{vyper}, to write smart contracts
that compile to the EVM bytecode. 



\noindent \textbf{Decentralized Finance (DeFi):}
DeFi is a peer-to-peer financial ecosystem built on top of
blockchains~\cite{wust2018you}. The building blocks of DeFi are smart contracts
that manage digital assets. 
A few DeFi protocols dominate the DeFi market and serve as references for other
decentralized applications: stable coins (e.g., USDC and USDT), price oracles,
decentralized exchanges, and lending and borrowing platforms. 
In DeFi, a special type of loan called \textbf{flash
loan} allows lenders to offer loans to borrowers without upfront collaterals
deposits. The loan is only valid within a single transaction and must be repaid
with fees before the completion of the transaction.
 


%\subsection{Transaction Execution} 

%A transaction is constituted of a sender address, a recipient address, a
%transferred value of native token (can be zero), a transaction data (can be
%empty), and a gas value to pay the transaction fees. If the transaction
%recipient address is associated with a user account then the transaction is
%called \emph{simple payment} transaction that transfers the value of native
%token from the sender account to the recipient account and the transaction's
%data is empty in this case. Otherwise, if the transaction recipient address is
%associated with a smart contract account then the transaction's data identifies
%a method of the recipient smart contract's code together with arguments passed
%to execute the method. When the transaction is received by the block miners,
%the corresponding smart contract's method is executed in the EVM, modifying the
%smart contract's balance and storage trie accordingly. Note that the execution
%of each EVM instruction, such as read/write operations on the underlying
%storage trie, is associated with a gas fee. The transaction's gas value must
%exceed the accumulated gas fees at the end of execution, otherwise the
%transaction is aborted~\cite{ethereum22,buterin2014ethereum}.



